% 빌드: bash scripts/build_paper.sh  (XeLaTeX + BibTeX)
\documentclass[11pt,a4paper]{article}

\usepackage{fontspec}
\IfFontExistsTF{Noto Serif CJK KR}{%
  % 한글에는 이탤릭이 없으므로 \emph는 고딕으로 표시한다
  \setmainfont{Noto Serif CJK KR}[BoldFont={* Bold}, ItalicFont={Noto Sans CJK KR}, BoldItalicFont={Noto Sans CJK KR Bold}]
  \setsansfont{Noto Sans CJK KR}[BoldFont={* Bold}]
}{%
  \IfFontExistsTF{AppleMyungjo}{\setmainfont{AppleMyungjo}\setsansfont{Apple SD Gothic Neo}}%
  {\setmainfont{Batang}\setsansfont{Malgun Gothic}}%
}
\XeTeXlinebreaklocale "ko"
\XeTeXlinebreakskip = 0pt plus 1pt

\usepackage[margin=25mm]{geometry}
\usepackage{setspace}\setstretch{1.35}
\usepackage{graphicx,booktabs,amsmath,xcolor,caption}
\usepackage[round,authoryear]{natbib}
\usepackage[hidelinks]{hyperref}
\captionsetup{font=small,labelfont=bf,labelsep=period}
\renewcommand{\figurename}{그림}
\renewcommand{\tablename}{표}
\renewcommand{\refname}{참고문헌}
\renewcommand{\abstractname}{초록}
\renewcommand{\contentsname}{차례}
\renewcommand{\appendixname}{부록}
\setlength{\parindent}{1em}

% 미완성 표시. scripts/check_submission.py는 이 표시가 남아 있으면 제출 묶기를 거부한다.
\newcommand{\TODO}[1]{\textcolor{red}{\textbf{[TODO: #1]}}}

% 표지 정보. 작성자 이름은 제출 전에 실명으로 바꿔도 된다.
\newcommand{\PaperTitle}{미리 주는 화면 테두리 신호는 반응을 빠르게 하는가:\\[2pt]
  \large 개인 기기에서 수행한 웹 기반 단순 반응시간 실험}
\newcommand{\PaperAuthor}{rnjswnghl\\[2pt]\small\url{https://github.com/rnjswnghl/reflex-hub}}
\newcommand{\PaperDate}{작성일: \today}
\newcommand{\ExperimentURL}{\url{https://reflex-hub.vercel.app}}

% 결과 수치는 analysis/run_analysis.py가 generated/results.tex로 만든다. 없으면 '대기' 값을 쓴다.
\IfFileExists{generated/results.tex}{\input{generated/results}\def\ResPending{0}}{% 데이터 수집 전 기본값. analysis/run_analysis.py를 실행하면 generated/results.tex가 이 파일 대신 쓰인다.
% 매크로 목록은 run_analysis.py의 write_tex()와 같아야 한다 (scripts/check_submission.py가 대조).
\def\ResPending{1}
\newcommand{\Pend}{\textcolor{red}{[대기]}}
\newcommand{\ResSimulated}{0}
\newcommand{\Nrecruited}{\Pend}
\newcommand{\Nexcluded}{\Pend}
\newcommand{\NexclPilot}{\Pend}
\newcommand{\NexclWindow}{\Pend}
\newcommand{\NexclRepeat}{\Pend}
\newcommand{\NexclIncomplete}{\Pend}
\newcommand{\NexclFewValid}{\Pend}
\newcommand{\Nanalyzed}{\Pend}
\newcommand{\TrialsAnalyzed}{\Pend}
\newcommand{\TrialsExcluded}{\Pend}
\newcommand{\TrialsExcludedPct}{\Pend}
\newcommand{\CollectStart}{\Pend}
\newcommand{\CollectEnd}{\Pend}
\newcommand{\NinputTouch}{\Pend}
\newcommand{\NinputMouse}{\Pend}
\newcommand{\NinputKey}{\Pend}
\newcommand{\NinputMixed}{\Pend}
\newcommand{\MedCue}{\Pend}
\newcommand{\MedNocue}{\Pend}
\newcommand{\HL}{\Pend}
\newcommand{\HLCIlo}{\Pend}
\newcommand{\HLCIhi}{\Pend}
\newcommand{\HLupper}{\Pend}
\newcommand{\Wplus}{\Pend}
\newcommand{\Wp}{\Pend}
\newcommand{\DiffMedian}{\Pend}
\newcommand{\DiffMean}{\Pend}
\newcommand{\DiffSD}{\Pend}
\newcommand{\DiffCIlo}{\Pend}
\newcommand{\DiffCIhi}{\Pend}
\newcommand{\DiffUpper}{\Pend}
\newcommand{\Tstat}{\Pend}
\newcommand{\Tdf}{\Pend}
\newcommand{\Tp}{\Pend}
\newcommand{\Dz}{\Pend}
\newcommand{\Kfaster}{\Pend}
\newcommand{\Kslower}{\Pend}
\newcommand{\Nnonzero}{\Pend}
\newcommand{\SignP}{\Pend}
\newcommand{\Imbalance}{\Pend}
\newcommand{\MDEdz}{\Pend}
\newcommand{\SensFloorHL}{\Pend}
\newcommand{\SensFloorP}{\Pend}
\newcommand{\SensFloorN}{\Pend}
\newcommand{\SensTrimHL}{\Pend}
\newcommand{\SensTrimP}{\Pend}
\newcommand{\FAcue}{\Pend}
\newcommand{\FAnocue}{\Pend}
\newcommand{\FAcuePct}{\Pend}
\newcommand{\FAnocuePct}{\Pend}
\newcommand{\CatchPerCond}{\Pend}
\newcommand{\FAdiffP}{\Pend}
\newcommand{\AntCue}{\Pend}
\newcommand{\AntNocue}{\Pend}
\newcommand{\FastCuePct}{\Pend}
\newcommand{\FastNocuePct}{\Pend}
\newcommand{\SawYes}{\Pend}
\newcommand{\SawNo}{\Pend}
\newcommand{\SawUnsure}{\Pend}
\newcommand{\LagMin}{\Pend}
\newcommand{\LagMax}{\Pend}
\newcommand{\CueOnMin}{\Pend}
\newcommand{\CueOnMax}{\Pend}
\newcommand{\TrialsPerCond}{\Pend}
\newcommand{\ResVerdict}{\Pend}
\newcommand{\ResSupported}{0}
}
\providecommand{\Pend}{\textcolor{red}{[대기]}}

\title{\PaperTitle}
\author{\PaperAuthor}
\date{\PaperDate}

\begin{document}
\maketitle

\ifnum\ResPending=1
  \begin{center}\colorbox{red!10}{\parbox{0.9\linewidth}{\centering\color{red!70!black}
  \TODO{데이터 수집 전 원고입니다. 결과 수치와 그림은 분석 실행 후 자동으로 채워집니다.}}}\end{center}
\fi
\ifnum\ResSimulated=1
  \begin{center}\colorbox{red!10}{\parbox{0.9\linewidth}{\centering\color{red!70!black}\bfseries
  \TODO{가상(SIMULATED) 데이터로 만든 점검용 빌드입니다. 제출 불가.}}}\end{center}
\fi

% 초록은 본문이 끝난 뒤 각 절의 첫 문장을 이어 붙여 쓴다(과제 안내 5/5).
% 수치는 모두 generated/results.tex 매크로에서 온다. 숫자를 손으로 적지 않는다.
\begin{abstract}
\noindent
\textbf{문제의식.} 게임과 앱은 ``곧 무언가 나온다''는 시각 효과를 자주 쓰지만, 그런 효과가 사용자의 반응을 실제로 빠르게 하는지는 대개 확인 없이 쓰인다.
실험실 연구에서는 목표 직전의 중립적 경고 신호가 단순 반응시간을 줄인다고 알려져 있다.
이 연구는 그 결과가 참가자 각자의 휴대폰·PC에서, 감독 없이 하는 짧은 웹 과제에서도 나타나는지 확인했다.
\textbf{방법.} 직접 만든 웹 실험 페이지에서 참가자는 화면 가운데 원이 나타나면 최대한 빨리 눌렀다.
절반의 시행에서는 원이 나오기 500\,ms 전에 화면 테두리가 100\,ms 동안 밝아졌고(경고 신호 있음), 나머지 절반에서는 아무 변화가 없었다.
두 조건은 한 참가자 안에서 시행마다 무작위로 섞였다(1인당 본 시행 40회).
분석 규칙은 데이터 수집 전에 고정해 공개 저장소에 기록했다.
\textbf{결과.} \CollectStart{}부터 \CollectEnd{}까지 \Nrecruited{}명이 참여했고, 사전에 정한 제외 기준을 적용한 뒤 \Nanalyzed{}명을 분석했다.
참가자별 반응시간 중앙값의 차이(신호 있음 $-$ 없음)는 중앙값 \DiffMedian{}\,ms(부트스트랩 95\% CI [\BootCIlo, \BootCIhi])였고,
사전 등록한 단측 Wilcoxon 부호순위 검정은 $p$ \Wp{}였다.
\Nnonzero{}명 중 \Kfaster{}명이 신호가 있을 때 더 빨랐다.
\textbf{결론.} 사전에 정한 판정 규칙에 따라 가설은 \ResVerdict.
\ifnum\ResSupported=1
화면 테두리 점등처럼 UI 효과에 가까운 신호도, 통제되지 않은 개인 기기 환경에서 같은 사람의 반응을 빠르게 했다.
다만 이 설계는 신호가 주는 각성과 ``언제 나올지''에 대한 시간 정보를 분리하지 못한다.
\else
이 결과는 ``효과가 없다''가 아니라 이 표본과 환경에서 효과를 확인하지 못했다는 뜻이며, 이 표본이 잡아낼 수 있는 최소 효과크기는 $d_z \approx$ \MDEdz{}였다.
\fi

\medskip\noindent
\textbf{주제어:} 경고 신호, 각성, 단순 반응시간, 웹 실험, 참가자 내 설계, 사전 등록
\end{abstract}

\section{서론}

\subsection{문제의식}
리듬 게임의 판정선 깜빡임, 경매 앱의 마감 직전 테두리 강조, 퀴즈 앱의 ``준비'' 애니메이션처럼, 인터페이스는 중요한 순간 직전에 시각 신호를 자주 보낸다.
이런 효과는 사용자가 더 빨리 반응하게 돕는다고 기대하고 쓰지만, 그 기대가 실제 사용 환경에서 맞는지는 대개 확인되지 않는다.
이 연구는 이 질문을 검증 가능한 한 문장으로 좁혀, 직접 만든 웹 실험으로 데이터를 모아 확인한다.

\subsection{이미 알려진 것}
주의 연구는 주의를 하나의 기능이 아니라 여러 하위 체계로 본다.
\citet{posner1990}은 경고 신호에 반응해 준비 상태를 높이는 \emph{각성}(alerting)을 정향, 실행 통제와 구분되는 기능으로 제안했고,
\citet{petersen2012}은 이후 20년의 연구를 바탕으로 이 구분을 갱신했다.
중립적 경고 신호, 즉 목표가 어디에 나올지나 무엇인지는 알려 주지 않는 신호 뒤에는 신호가 없을 때보다 반응이 빠르다 \citep{weinbach2012}.
\citet{fan2002}은 경고 단서가 있는 시행과 없는 시행의 반응시간 차이로 각성 효과를 재는 과제를 만들었다.
이 과제는 성인 40명에게서 개인 단위로도 신뢰할 만한 추정치를 냈다.
경고와 목표 사이 간격(예비 기간)이 단순 반응시간에 영향을 준다는 사실도 오래전부터 정리되어 있다 \citep{niemi1981}.
이 때문에 이 연구는 두 조건의 대기시간 분포를 똑같이 맞췄다(\ref{sec:design}절).

\subsection{웹으로 반응시간을 재도 되는가}
반응시간 연구는 전통적으로 실험실의 전용 장비에서 이뤄졌다.
전용 측정 장비로 웹 실험을 점검한 연구들은, 브라우저에서 측정한 반응시간이 실험실보다 변동이 조금 크지만
참가자 내 조건 비교에는 충분히 쓸 수 있다고 보고했다 \citep{reimers2015,bridges2020,anwylirvine2021}.
특히 \citet{pronk2020}은 터치스크린과 키보드 기기에서 반응시간이 일관되게 과대 측정되며 그 크기가 기기마다 다르다고 보고했다.
같은 연구는 이런 기기 간 차이가 절대 반응시간보다 같은 사람 안의 상대 반응시간에 덜 영향을 준다고 정리했다.
이 결과가 이 연구가 참가자 간 비교가 아니라 \emph{참가자 내} 설계를 택한 핵심 근거다.

\subsection{이 연구가 새로 확인하려는 것}
기존 연구가 보인 것은 (1) 통제된 실험실 환경에서 경고 신호가 반응을 빠르게 한다는 것과
(2) 웹 브라우저가 조건 간 반응시간 차이를 잡아낼 만큼 정밀하다는 것이다.
이 연구는 두 결과를 합쳐, 참가자가 \emph{자기 기기로, 감독 없는 환경에서} 하는 3분짜리 웹 과제에서도,
UI 효과처럼 생긴 \emph{화면 테두리 점등}이 같은 사람의 반응시간을 줄이는지 확인한다.
새로운 이론을 세우는 연구가 아니라, 알려진 효과를 현실적인 사용 환경에서 다시 확인하는 개념적 재현이다.

\subsection{가설}
\begin{quote}
\textbf{가설.} 목표 자극이 나타나기 500\,ms 전에 화면 테두리가 100\,ms 동안 밝아지는 시각 경고 신호가 있으면,
신호가 없을 때보다 같은 참가자의 단순 반응시간이 짧아진다.
\end{quote}
원인은 경고 신호의 유무이고, 결과는 단순 반응시간(ms)이며, 방향은 ``짧아진다''이다.
가설이 맞다면 대다수 참가자에게서 [신호 있음 중앙값 $-$ 신호 없음 중앙값]이 음수로 나오고,
사전 등록한 검정에서 그 중앙값이 0보다 작다는 결과가 나와야 한다.
가설이 틀렸다면 이 차이가 0 근처에 흩어지거나, 경고가 오히려 방해가 되어 양수가 많이 나온다.

\section{방법}

\subsection{설계}\label{sec:design}
경고 신호(있음/없음)를 참가자 내 요인으로 둔 단일 요인 설계다.
한 참가자가 두 조건을 모두 수행하고, 조건 순서는 참가자마다 시행 단위로 무작위 배정했다.
이렇게 하면 기기의 고정 입력 지연과 개인의 기본 반응속도가 두 조건에 똑같이 들어가, 조건 간 차이에서 상쇄된다 \citep{pronk2020}.
표~\ref{tab:design}에 바꾼 조건과 고정한 조건을 정리했다.

\begin{table}[h]
\centering\small
\caption{실험 설계: 바꾼 조건과 고정한 조건}\label{tab:design}
\begin{tabular}{p{0.24\linewidth}p{0.68\linewidth}}
\toprule
\multicolumn{2}{l}{\textbf{바꾼 조건(독립변수)}}\\
신호 있음 & 목표 500\,ms 전, 화면 안쪽 테두리(두께 10\,px, 밝은 회백색)가 100\,ms 점등. 1인당 20회\\
신호 없음 & 같은 시점에 화면 변화 없음. 1인당 20회\\
\midrule
\multicolumn{2}{l}{\textbf{고정한 조건(통제변수)}}\\
경고 시점 $\to$ 목표 & 두 조건 모두 500\,ms (신호 없음 조건에서는 보이지 않는 가상 시점)\\
시행 시작 $\to$ 경고 시점 & 1000--2500\,ms 균등분포에서 시행마다 추출, 두 조건 같은 분포\\
목표 자극 & 화면 중앙 흰 원, 지름 120\,px (CSS 픽셀)\\
응답 & 화면 아무 곳 터치·클릭(\texttt{pointerdown}), 또는 Space/Enter 키(\texttt{keydown})\\
반응 창 & 목표 후 1500\,ms. 넘으면 놓침으로 기록\\
시행 수 & 연습 6회(조건별 3회, 분석 제외) + 본 시행 40회\\
\bottomrule
\end{tabular}
\end{table}

\subsection{참가자}
연구자의 지인과, 연구자가 실험 링크를 올린 온라인 커뮤니티의 이용자를 대상으로 편의 표집했다.
수집 기간은 링크를 처음 공유한 날부터 7일로 미리 정했다.
실제로는 \CollectStart{}부터 \CollectEnd{}까지 모았고, 이 기간에 본 실험을 끝낸 세션은 \Nrecruited{}개였다.
목표 인원은 20명 이상으로 정했다.
이 숫자는 검정력 분석이 아니라 1주 안에 모을 수 있는 범위를 기준으로 정했으며, 그 한계는 \ref{sec:limits}절에 적었다.
이름·연락처 등 개인을 식별하는 정보는 받지 않았다.
연령대는 선택 응답으로 받았고, 참가자는 시작 화면에서 기록되는 항목을 확인한 뒤 동의했다.

\subsection{장치와 자극}
실험 페이지는 React와 Vite로 만들었고 Vercel에 배포했다(\ExperimentURL).
참가자는 각자의 휴대폰, 태블릿, PC 브라우저로 참여했다.
자극 제시와 시간 측정은 React 렌더링을 거치지 않고 DOM을 직접 바꾸는 별도 모듈에서 처리했다.
목표 자극은 정해진 시각에 가장 가까운 \texttt{requestAnimationFrame} 콜백에서 그렸고, 그 시점의 \texttt{performance.now()}를 제시 시각으로 기록했다.
반응 시각은 입력 이벤트의 \texttt{event.timeStamp}로 기록했다.
\texttt{click}은 손을 뗄 때 발생해 불필요한 지연이 섞이므로, 누르는 순간 발생하는 \texttt{pointerdown}을 썼다.
난수는 \texttt{crypto.getRandomValues}로 만들었다.

\subsection{절차}
참가자는 안내와 동의 화면 다음에 방법 설명을 읽고 연습 6회를 수행했다.
매 시행은 화면 가운데 작은 점이 나타나며 시작된다.
1000--2500\,ms 뒤의 경고 시점에, 신호 있음 조건이면 테두리가 100\,ms 동안 밝아졌다.
그로부터 500\,ms 뒤 점이 사라지고 흰 원이 나타났다.
참가자는 원이 보이면 최대한 빨리 화면을 누르거나 키를 눌렀다.
원이 나오기 전에 누르면 ``너무 일러요''라는 안내와 함께 그 시행을 조기 반응으로 기록했다.
반응 뒤에는 그 시행의 반응시간을 600\,ms 동안 보여 주고, 500\,ms 쉰 뒤 다음 시행으로 넘어갔다.
본 시행 40회를 마치면 개인 결과를 보여 주고 기록을 서버로 보냈다.
전체 소요 시간은 약 3분이었다.

\subsection{측정값}
종속변수는 단순 반응시간, 즉 반응 시각에서 목표 제시 시각을 뺀 값(ms)이다.
시행마다 조건, 계획한 대기시간, 실제 경고--목표 간격, 상태(정상·조기 반응·놓침), 반응시간, 입력 방식, 시행 도중 탭 전환 여부를 기록했다.
세션마다 앱 버전, 같은 브라우저의 재참여 여부, 포인터 종류, 화면 크기, 추정 프레임 간격을 기록했다.
참가자 1명의 기록은 JSON 파일 하나로 저장했다.
모든 원자료와 정리된 CSV는 재현 패키지에 들어 있다.

\subsection{측정 도구 점검}
수집 전에 브라우저 안의 ``로봇 참가자''로 페이지를 점검했다.
로봇은 원이 뜨면 정확히 250\,ms 뒤에 누르도록 만들었다.
데스크톱 Chromium에서 46행(연습 6, 본 시행 40)이 설계대로 기록되었다.
실제 경고--목표 간격은 499.4--500.3\,ms(설계 500\,ms)였고, 기록된 반응시간은 250.2--250.5\,ms였다.
일부러 넣은 조기 반응 1건은 반응시간 없이 조기 반응으로 기록되었다.
즉 소프트웨어 측정 오차는 1\,ms 미만이었다.
다만 로봇은 이벤트를 소프트웨어로 발생시키므로, 이 점검에는 터치스크린·키보드 하드웨어의 지연이 포함되지 않는다.

\subsection{사전 등록한 분석}
분석 규칙은 데이터 수집 전에 문서와 코드로 고정했고, 공개 Git 저장소의 커밋 기록으로 시점을 남겼다.
\textbf{시행 제외}: 연습 시행, 조기 반응, 놓침, 시행 도중 탭·앱 전환이 감지된 시행, 반응시간 100\,ms 미만.
\textbf{참가자 제외}(순서대로 적용):
같은 브라우저의 두 번째 이후 참여,
본 시행을 모두 마치지 않은 세션,
유효 시행 30개 미만이거나 어느 한 조건의 유효 시행이 10개 미만인 경우.
\textbf{분석 단위}: 참가자별·조건별 유효 반응시간의 중앙값과, 그 차이 $d = $ 중앙값(신호 있음) $-$ 중앙값(신호 없음).
반응시간 분포는 오른쪽 꼬리가 길어, 평균 대신 이상값의 영향을 덜 받는 중앙값을 썼다 \citep{ratcliff1993}.
\textbf{확증 검정}: $d$에 대한 단측 Wilcoxon 부호순위 검정(대립가설: 중앙값 $<0$), $\alpha = .05$.
$p<.05$이고 $d$의 중앙값이 음수이면 ``가설이 지지되었다'', 그 외에는 ``지지되지 않았다''로 판정한다.
\textbf{보조 분석}: 단측 대응 $t$검정, 효과크기 $d_z$, 평균 차이의 95\% 신뢰구간,
$d$ 중앙값의 부트스트랩 95\% 신뢰구간(재표집 10{,}000회, 시드 20261030), 부호 검정, 민감도 분석.
\textbf{탐색 분석}(확증 아님): 입력 방식별 $d$, 조건별 조기 반응 건수.
분석은 Python(pandas, SciPy)으로 했고, 결과 수치와 그림은 분석 스크립트가 원자료에서 바로 만들어 이 문서에 넣었다.

\section{결과}
% 이 절의 수치는 모두 analysis/run_analysis.py가 만든 매크로다. 숫자를 손으로 고치지 않는다.

\subsection{참가자와 시행}
수집 기간에 서버에 도착한 완료 세션은 \Nrecruited{}개였다.
사전 등록한 기준으로 \Nexcluded{}개를 제외했다.
사유는 파일럿 \NexclPilot{}개, 수집 기간 밖 \NexclWindow{}개, 재참여 \NexclRepeat{}개, 미완료 \NexclIncomplete{}개, 유효 시행 부족 \NexclFewValid{}개였다.
남은 \Nanalyzed{}명의 목표 시행 \TrialsAnalyzed{}회 중 \TrialsExcluded{}회(\TrialsExcludedPct\%)가 유효 시행 정의를 만족하지 못해 빠졌다.
참가자별 유효 시행 수의 조건 간 차이는 평균 \Imbalance{}회였다.
분석한 참가자의 주 입력 방식은 터치 \NinputTouch{}명, 마우스 \NinputMouse{}명, 키보드 \NinputKey{}명, 혼합 \NinputMixed{}명이었다.
기록된 경고--목표 간격은 \LagMin--\LagMax\,ms, 테두리 점등 시간은 \CueOnMin--\CueOnMax\,ms였다.
조작 점검에서 테두리를 봤다는 응답은 \SawYes{}명, 못 봤다 \SawNo{}명, 모르겠다 \SawUnsure{}명이었다.
제외된 세션과 시행의 사유는 재현 패키지의 \texttt{analysis/output/exclusions.csv}와 \texttt{results.json}에 있다.

\subsection{기술통계}
표~\ref{tab:desc}는 참가자별 조건 중앙값을 다시 참가자들에 걸쳐 요약한 것이다.
그림~\ref{fig:paired}은 참가자 한 명을 선 하나로 이어, 두 조건의 값을 나란히 보여 준다.

\begin{table}[h]
\centering\small
\caption{참가자별 반응시간 중앙값의 요약 (단위 ms, $n=\Nanalyzed$)}\label{tab:desc}
\IfFileExists{generated/table_desc.tex}{\input{generated/table_desc}}{\Pend}
\end{table}

\begin{figure}[h]
\centering
\IfFileExists{figures/fig_paired.pdf}{\includegraphics[width=0.48\linewidth]{figures/fig_paired.pdf}}{\fbox{\Pend}}\hfill
\IfFileExists{figures/fig_diff.pdf}{\includegraphics[width=0.48\linewidth]{figures/fig_diff.pdf}}{\fbox{\Pend}}
\caption{왼쪽: 참가자별 조건 중앙값. 회색 선 하나가 참가자 한 명이고, 검은 선은 참가자 중앙값이다.
오른쪽: 참가자별 차이(신호 있음 $-$ 없음). 점 하나가 참가자 한 명이고, 세로선은 Hodges--Lehmann 추정치, 음영은 그 95\% 신뢰구간이다.
0보다 왼쪽이면 신호가 있을 때 빨랐다는 뜻이다.}\label{fig:paired}
\end{figure}

\subsection{확증 분석}
참가자별 차이 $d$의 Hodges--Lehmann 추정치는 \HL\,ms(95\% CI [\HLCIlo, \HLCIhi], 단측 95\% 상한 \HLupper)였다.
사전 등록한 단측 Wilcoxon 부호순위 검정 결과는 $W^{+}=\Wplus$, $p$ \Wp{}였다.
사전에 정한 판정 규칙에 따라 가설은 \textbf{\ResVerdict}.
참고로 $d$의 표본 중앙값은 \DiffMedian\,ms였다.
이 값은 위에 적은 조건별 참가자 중앙값(신호 있음 \MedCue\,ms, 없음 \MedNocue\,ms)의 차와 같지 않다.
중앙값의 차는 차의 중앙값과 일반적으로 다르기 때문이다.

\subsection{보조 분석과 민감도 분석}
평균 차이는 \DiffMean\,ms(표준편차 \DiffSD, 95\% CI [\DiffCIlo, \DiffCIhi], 단측 95\% 상한 \DiffUpper)였다.
단측 대응 $t$검정은 $t(\Tdf)=\Tstat$, $p$ \Tp{}였고, 효과크기는 $d_z=\Dz$였다.
차이가 0이 아닌 \Nnonzero{}명 중 \Kfaster{}명이 신호가 있을 때 더 빨랐고, \Kslower{}명은 더 느렸다(부호 검정 $p$ \SignP).
반응시간 하한을 150\,ms로 올리면($n=\SensFloorN$) HL 추정치는 \SensFloorHL\,ms, $p$ \SensFloorP{}였다.
중앙값 대신 10\% 절사평균을 쓰면 HL 추정치는 \SensTrimHL\,ms, $p$ \SensTrimP{}였다.
이 표본 크기에서 단측 대응 $t$검정이 검정력 .80으로 잡아낼 수 있는 최소 효과크기는 $d_z\approx\MDEdz$이다(주 검정인 Wilcoxon 기준이 아니다).

\subsection{전략 점검}
캐치 시행 오반응은 신호 있음 조건에서 \FAcue{}건(\FAcuePct\%), 신호 없음 조건에서 \FAnocue{}건(\FAnocuePct\%)이었다.
참가자별 오반응 수의 조건 차이에 대한 양측 부호순위 검정은 $p$ \FAdiffP{}였다.
조기 반응은 신호 있음 \AntCue{}건, 신호 없음 \AntNocue{}건이었다(조건별 목표 시행 \TrialsPerCond{}회 중).
100--150\,ms 반응의 비율은 신호 있음 \FastCuePct\%, 신호 없음 \FastNocuePct\%였다.

\subsection{탐색 분석 (사전 등록한 확증 분석이 아님)}
입력 방식별, 조작 점검 응답별 차이는 \texttt{results.json}의 \texttt{exploratory}에 있다.
집단별 인원이 적어 해석하지 않는다.

\TODO{데이터를 본 뒤 작성: 가설대로 나온 부분과 다르게 나온 부분을 나눠 한 문단으로 정리 (T10-C32, C35). 반대 방향 참가자 수(\Kslower)와 그 크기, 전략 점검이 가리키는 바를 숨기지 말고 적는다.}

\section{논의}

\subsection{결과 요약}
이 연구는 화면 테두리를 잠깐 밝히는 시각 경고 신호가, 개인 기기에서 하는 웹 과제에서도 같은 사람의 단순 반응시간을 줄이는지 확인했다.
사전에 정한 규칙에 따라 가설은 \ResVerdict.

% 아래 두 갈래의 해석 틀은 데이터 수집 전에 써 두었다(Git 기록 참고). 결과가 나온 뒤 고르는 것이 아니라 미리 정한 대로 따른다.
\ifnum\ResSupported=1
\subsection{해석}
결과는 실험실 연구의 각성 효과 \citep{weinbach2012,fan2002}와 같은 방향이다.
이 연구의 기여는 효과의 존재가 아니라 환경에 있다.
참가자는 서로 다른 기기와 화면 크기, 입력 방식으로, 감독 없이 실험에 참여했다.
그런 조건에서도 UI 장식에 가까운 테두리 점등이 반응을 빠르게 했다.
따라서 ``곧 나온다''는 시각 신호는 실제 사용 환경에서도 반응 준비를 돕는다고 볼 수 있다.
\else
\subsection{해석}
이 결과를 ``경고 신호는 효과가 없다''로 읽으면 안 된다.
가능한 설명은 미리 네 가지로 정해 두었고, 원자료로 확인할 수 있는 것부터 확인한다.
(1) \emph{표본 부족}: 이 표본이 잡아낼 수 있는 최소 효과크기($d_z\approx\MDEdz$)보다 실제 효과가 작을 수 있다.
(2) \emph{신호를 못 봄}: 테두리는 화면 가장자리에 있어, 휴대폰을 쥔 손이나 시선 위치에 따라 눈에 덜 띌 수 있다. 입력 방식별 차이로 일부 확인한다.
(3) \emph{시간 정보 효과의 상쇄}: 신호 없음 조건에서도 기다릴수록 목표가 곧 나온다는 예측이 생겨 반응이 빨라진다 \citep{niemi1981}. 이 효과가 두 조건의 차이를 줄였을 수 있다.
(4) \emph{성급한 반응과의 교환}: 신호가 반응을 앞당기는 대신 조기 반응을 늘렸다면, 그 시행은 분석에서 빠져 효과가 작게 보인다. 조건별 조기 반응 건수로 확인한다.
\fi

\TODO{데이터를 본 뒤 작성: 위 설명 가운데 원자료로 확인한 것과 그 결과(T10-C36). 가설을 고쳤다면 고치기 전 문장과 고친 문장, 고친 이유(T10-C37). 고치지 않았다면 그 이유.}

\subsection{한계}\label{sec:limits}
\begin{enumerate}
\item \textbf{메커니즘을 분리하지 못한다.} 경고 신호는 각성을 높이는 동시에 목표가 정확히 500\,ms 뒤에 나온다는 시간 정보도 준다.
이 두 과정은 대부분의 설계에서 섞여 있고 \citep{weinbach2012}, 이 실험도 둘을 구분하지 못한다.
따라서 이 연구가 말할 수 있는 것은 ``이 신호가 반응시간을 바꿨는가''까지이며, 그 이유가 각성인지 시간 예측인지는 말할 수 없다.
\item \textbf{하드웨어 지연은 측정하지 않았다.} 기기마다 반응시간이 수십 ms 이상 과대 측정될 수 있다 \citep{pronk2020}.
참가자 내 차이는 이 고정 지연을 상쇄하지만, 이 논문의 \emph{절대} 반응시간을 다른 연구와 직접 비교해서는 안 된다.
\item \textbf{표본.} 지인과 커뮤니티 이용자의 편의 표본이고, 인원은 검정력 분석이 아니라 수집 가능 범위로 정했다. 결과를 일반 인구로 넓히기 어렵다.
\item \textbf{감독 없는 환경.} 참가자의 주의 상태, 주변 소음, 화면 밝기를 통제하지 못했다. 탭 전환은 감지해 제외했지만 화면 밖의 방해는 감지할 수 없다.
\item \textbf{재참여 판별.} 재참여는 같은 브라우저의 저장소로만 판별하므로, 다른 기기나 시크릿 창으로 다시 참여한 경우는 걸러내지 못한다.
\item \textbf{자극 크기.} 테두리 두께와 원 크기는 CSS 픽셀로 고정했다. 화면 크기에 따라 실제 시각 크기와 눈에 띄는 정도가 달랐다.
\end{enumerate}

\subsection{결론}
\TODO{데이터를 본 뒤 작성: 한 문단 이상(T10-C38). 원자료와 위 분석에 있는 내용만 쓴다(T10-C40).}


\bibliographystyle{apalike}
\bibliography{refs}

\appendix
\section{재현 방법}\label{app:repro}
재현 패키지(ZIP)를 풀거나 \url{https://github.com/rnjswnghl/reflex-hub}를 받은 뒤 아래를 실행하면, 원자료에서 이 논문의 모든 수치와 그림이 다시 만들어진다.
자세한 내용은 패키지의 \texttt{REPRODUCE.md}에 있다.

\begin{enumerate}
\item 환경: Python 3.10 이상, \texttt{pip install -r analysis/requirements.txt}. 논문 빌드에는 XeLaTeX이 필요하다.
\item 원자료 정리: \texttt{python3 analysis/build\_dataset.py}\\
\texttt{data/raw/sessions/*.json} $\to$ \texttt{data/raw/trials.csv}, \texttt{sessions.csv}
\item 사전 등록한 분석: \texttt{python3 analysis/run\_analysis.py}\\
$\to$ \texttt{analysis/output/results.json}, 논문 수치 매크로, 그림
\item 논문 빌드: \texttt{bash scripts/build\_paper.sh}
\item 실험 페이지 직접 실행: \texttt{cd web \&\& npm install \&\& npm run dev}
\end{enumerate}

\section{AI에게 맡긴 일과 직접 판단한 일}\label{app:ai}
이 연구는 AI(Claude)와 함께 진행했다.
역할 구분과 판단 근거는 재현 패키지의 \texttt{docs/06-ai-and-me.md}에 기록했고, 요약은 다음과 같다.
AI는 가설 문장의 결함 지적, 실험 페이지와 분석 코드 작성, 문헌 후보 제시와 서지 대조, 원고 초안 작성을 맡았다.
AI가 제시한 문헌은 모두 웹에서 직접 서지를 대조했고, 대조 과정에서 원문과 맞지 않는 연결 문장 1건을 고쳤다.
분석 규칙은 데이터를 보기 전에 고정했으므로, AI도 연구자도 결과를 본 뒤 판정 기준을 바꿀 수 없었다.

\section{문서 대응표}
과제 통과 기준(T10-C01--C51)과 이 논문·패키지의 대응 위치는 \texttt{docs/CHECKLIST.md}에 있다.


\end{document}
